\section{Experiment Design}

\subsection{Research object, channels, and variables}

The ExperimentDesign Agent selects a computational-digital template. The scientific object is not a physical collider operation but a finite-resolution inference over scattering observables. The baseline hypothesis is a pointlike-SM prediction. Alternative families include specified dimension-six and dimension-eight contact operators, finite form-factor deformations, and excited-fermion benchmarks. Neutral- and charged-current Drell--Yan-like observables form the primary planned interfaces because their high-scale tails and angular structure are relevant to semileptonic interactions \cite{allwicher2023}. Lepton-hadron and future-facility channels are optional and are explicitly labeled as synthetic or projected until independently approved public inputs exist.

\begin{table*}[!t]
\caption{\texttt{DESIGN\_ONLY} modules of Resolution-Ladder Model Comparison.}
\label{tab:modules}
\centering
\footnotesize
\begin{tabularx}{\textwidth}{p{0.16\textwidth}p{0.24\textwidth}Y Y}
\toprule
\textbf{Module} & \textbf{Inputs and controls} & \textbf{Planned output} & \textbf{Failure or review condition}\\
\midrule
Channel contract & Momentum-transfer definition, binning, fiducial selection, status label, and response interface & Versioned channel manifest & Missing selection, incompatible $Q$ definition, or projection presented as data\\
Pointlike baseline & SM prediction order, PDFs, scale variations, backgrounds, and nuisance priors & Baseline expected yields and uncertainty envelope & Prediction provenance or nuisance correlations unavailable\\
Contact-EFT alternatives & Operator basis, coefficient set, flavor assumptions, cutoff, dimension-eight policy & EFT likelihood and validity diagnostic & Cutoff or truncation sensitivity changes the conclusion\\
Finite-size alternative & Declared form-factor expansion and range of use & Radius-like shape template & Expansion used outside justified kinematics or degenerate with response effects\\
Excited-state alternative & Mass, coupling, and compositeness-scale benchmark & Resonant or nonresonant template family & Benchmark not distinguishable from contact or background shape\\
Closure and holdout & Synthetic data, adversarial nuisance injection, training/held-out bins & Identifiability and predictive comparison & A nuisance model reproduces an apparent preference\\
\bottomrule
\end{tabularx}
\end{table*}

The independent variables are momentum transfer or invariant mass $Q$, angular variable $c$, EFT coefficient ratios $C_i/\Lambda^2$, declared dimension-eight terms $D_j/\Lambda^4$, form-factor parameters $r_f^2$, excited-state benchmark parameters, and nuisance vector $\eta$. The dependent objects are binned yields, normalized angular distributions, likelihood ratios, held-out predictive scores, and stability diagnostics. Controlled objects include binning version, theory order, PDF ensemble, fiducial definition, detector response, luminosity treatment, and data-release provenance.

\subsection{Generative model and likelihood}

For a channel $c$ and bin $b$, RLMC records an expected yield
\begin{equation}
\mu_{cb}(\theta,\eta)=
\mathcal{L}\int_{b}\left[R_{cb}\frac{d\sigma_c(\theta,\eta)}{d\Phi}\right]d\Phi
B_{cb}(\eta).
\label{eq:yield}
\end{equation}
Here $\mathcal{L}$ is integrated luminosity when a realized measurement is used, $R_{cb}$ is the response operator, $d\sigma_c/d\Phi$ is the differential prediction in phase space $\Phi$, $B_{cb}$ is the background contribution, $\theta$ collects physics parameters, and $\eta$ collects nuisance parameters. For synthetic closure tests, all terms are simulated inputs and are never described as observed data.

The joint likelihood is
\begin{equation}
\mathcal{L}(\theta,\eta)=
\prod_{c,b}\operatorname{Poisson}\left(n_{cb}\mid\mu_{cb}(\theta,\eta)\right)
\prod_k\pi_k(\eta_k),
\label{eq:likelihood}
\end{equation}
where $n_{cb}$ is a binned record only in an independently approved interface and $\pi_k$ represents declared nuisance constraints. The construction must preserve which nuisance parameters are correlated across channels and which are not. It is invalid to merge a future projection with a released measurement merely because both use the word sensitivity.

\subsection{EFT validity and cross-channel contract}

For a contact interaction, the inference includes a dimensionless diagnostic
\begin{equation}
\epsilon_{\mathrm{EFT}}=\frac{Q_{\max}}{\Lambda_{\mathrm{cut}}},
\label{eq:eftvalidity}
\end{equation}
with a predeclared acceptance policy rather than a universal numerical threshold. The analysis repeats the comparison over a cutoff sweep and an alternative operator-order policy. A point estimate at one cutoff is not sufficient evidence for compositeness when the conclusion is unstable under these choices.

Cross-channel consistency is tested only when the assumed model supplies a relation among coefficients or shape parameters. Channels may have distinct PDFs, response effects, and backgrounds, so RLMC shares a nuisance object only when physics and calibration provenance justify it. This avoids a false appearance of precision from forcing unrelated channels into one covariance matrix.

\section{Analysis and Decision Rules}

\subsection{Identifiability before preference}

The central question is not whether an alternative can fit a flexible curve. It is whether pointlike, contact, form-factor, and excited-state families have distinguishable predictions over the declared bins after nuisance variation. RLMC begins with synthetic closure tests generated from each family. It measures profile degeneracies, posterior correlations, and the ability to recover an injected class under held-out bins. If a finite form factor and a response-model deformation have indistinguishable effects, a radius interpretation is non-identifiable and must not be reported as a compositeness indication.

The primary comparison statistic is
\begin{equation}
\Delta T=-2\log
\left[
\frac{\mathcal{L}(\widehat{\theta}_{\mathrm{pointlike}},\widehat{\eta}_{\mathrm{pointlike}})}
{\mathcal{L}(\widehat{\theta}_{\mathrm{alternative}},\widehat{\eta}_{\mathrm{alternative}})}
\right].
\label{eq:teststat}
\end{equation}
Equation~\eqref{eq:teststat} does not establish coverage or model preference by itself. Its calibration must be checked with the appropriate pseudoexperiment, asymptotic, or Bayesian procedure selected by a qualified analysis team. The statistic is paired with held-out predictive performance, cutoff variation, and nuisance stress rather than interpreted in isolation.

\subsection{Stress tests and decision branches}

The first stress test varies the PDF ensemble and its correlations. The second perturbs detector response, background normalization, and luminosity terms within their declared uncertainty model. The third changes binning and withholds high-$Q$ bins. The fourth changes EFT order, operator subset, flavor assumption, and cutoff. The fifth injects controlled contact, form-factor, and excited-state signals into synthetic closure records. A conclusion is stable only if the relevant model remains identifiable and its improvement persists under predetermined stress tests.

\begin{table}[!t]
\caption{Predeclared RLMC decision branches.}
\label{tab:branches}
\centering
\footnotesize
\begin{tabularx}{\columnwidth}{p{0.27\columnwidth}Y}
\toprule
\textbf{Condition} & \textbf{Permitted conclusion}\\
\midrule
Identifiable, cutoff stable, nuisance robust, and cross-channel coherent alternative & The selected alternative is preferred within its declared model class and tested resolution range. It is not proof of a unique constituent.\\
Identifiable alternative with no predictive gain & The pointlike baseline remains adequate for tested observables and assumptions. Report a range-dependent constraint.\\
No hypothesis identifiability & No source-specific or constituent interpretation is warranted; revise bins, observables, or calibration.\\
EFT-order or cutoff instability & Do not quote a compositeness interpretation from that region; restrict or change the model.\\
Nuisance injection mimics deviation & Treat the result as a systematic-model ambiguity, not new physics.\\
\bottomrule
\end{tabularx}
\end{table}

\subsection{Reproducible reporting}

Every run produces a channel manifest, data-status label, theory calculator version, selection and binning record, nuisance ledger, EFT policy, synthetic closure seed, and decision log. The manifest is more than administrative metadata. It permits another analyst to identify whether a comparison used a real public measurement, a public likelihood recast, a synthetic record, or a future projection. This is essential when integrating SMEFTsim-like coefficient calculators \cite{brivio2017} with external experimental inputs.

\section{Expected Outcome Branches and Significance}

If RLMC finds that a pointlike-SM baseline is adequate under a validated public or future experimental interface, the correct result is a stronger and more transparent range-dependent constraint. It specifies the scales, channels, alternatives, and nuisance assumptions for which no additional substructure is required. Such a result does not close the possibility of deeper degrees of freedom beyond accessible resolution.

If an alternative is identifiable and remains preferred after the full stress program, the correct result is a model-conditioned anomaly or preference. It establishes neither a preon nor a final constituent. A contact coefficient can arise from many ultraviolet mechanisms; a form factor can have modeling ambiguities; an excited-state signal requires its own background and reconstruction tests. The next stage would be an explicit ultraviolet-model comparison and independent channel confirmation, not a statement that the smallest building block has been found.

If the analysis is non-identifiable, the negative result is useful. It identifies which nuisance or theory assumptions prevent a meaningful conclusion and motivates a better observable, a different kinematic regime, an improved response model, or a complementary facility. This is where high-energy, lepton-hadron, precision, and low-energy programs become strategically complementary \cite{abada2019,alekhin2016,jaeckel2010}.

\section{Limits and Human Review Requirements}

The proposal has substantive limits. PDFs are not optional details in proton-collision interpretations. A form-factor expansion has a finite domain of validity. EFT truncation is not eliminated by writing one operator in Eq.~\eqref{eq:eft}. An excited-state benchmark is not a complete compositeness theory. Cross-channel correlation is physically meaningful only when inputs and model justify it. A symmetry-respecting extension of the SM can parameterize possible departures, but parameterization is not evidence that a particular departure exists \cite{colladay1998}.

The workflow is also operationally limited. It does not operate accelerators, steer beams, configure triggers, control detectors, access restricted collider records, or publish a physics limit. Before any public-data application, a collider phenomenologist and experimental analyst must confirm data-release terms, fiducial definitions, detector response, likelihood validity, theory order, and statistical coverage. Before an external submission, a human must verify publisher or DOI landing pages because the requested application browser could not run in the Windows ACL sandbox.

\section{Conclusion}

The Standard Model currently provides the best-tested inventory of elementary matter fields: quarks and leptons, together with gauge and Higgs bosonic fields. Protons and neutrons are composite hadrons, not the final building blocks. The experimentally honest conclusion is that known quarks and leptons behave pointlike over tested scales and assumptions, not that every possible deeper layer has been ruled out forever.

RLMC turns this distinction into a concrete research program. It compares pointlike-SM, contact-EFT, finite-form-factor, and excited-fermion hypotheses across versioned channels, with an explicit nuisance ledger and EFT-validity gate. It treats a null result, a model preference, a systematic ambiguity, and non-identifiability as different scientific outcomes with different permitted conclusions. It asks not only what model is currently adequate, but also what evidence would genuinely require us to refine it.

\section{Resolution Road Map and Evidence Upgrading}

\subsection{A staged program rather than a single limit}

The question of pointlike behavior is best treated as a staged program. The first stage is synthetic closure. It asks whether the selected channel, binning, response model, and nuisance ledger can recover a known injected pointlike, contact, form-factor, or excited-state pattern. This is a test of the inference apparatus, not a stand-in for a collision measurement. It is especially valuable because a fitted contact coefficient can otherwise be driven by a poorly modeled tail or a hidden response correlation.

The second stage is channel qualification. A channel enters the common comparison only after its kinematic definition, fiducial selection, theory calculator, response model, and data status are registered. A public binned measurement without a response interface can be included as a qualitative evidence source but not as a quantitative likelihood. A public likelihood can be used only under its documented validity conditions. A future collider or electron-ion facility remains a projection until it produces released data. This staged admission protects the program from assigning equal weight to incomparable inputs.

The third stage is cross-channel synthesis. The combined analysis should not ask whether every channel has the same apparent deviation. Instead, it asks whether a particular physics model predicts a consistent coefficient, form-factor scale, or excited-state benchmark after the channels retain their own PDFs, backgrounds, and response effects. A disagreement may rule out the selected common model, expose an interface error, or point to a more complicated alternative. It does not become evidence for compositeness by itself.

\begin{table*}[!t]
\caption{Resolution-ladder evidence-upgrading rules.}
\label{tab:roadmap}
\centering
\footnotesize
\begin{tabularx}{\textwidth}{p{0.15\textwidth}p{0.23\textwidth}Y Y}
\toprule
\textbf{Stage} & \textbf{Admissible input} & \textbf{Primary question} & \textbf{Go or no-go rule}\\
\midrule
Closure & Synthetic records generated from stated hypotheses and nuisance models & Can the inference recover the injected class and separate it from adversarial nuisance shifts? & No-go for quantitative interpretation if classes are non-identifiable.\\
Qualification & Released binning, selections, response information, and theory provenance & Is the channel definition sufficient to construct a reproducible likelihood? & No-go for combined fit if status or response provenance is missing.\\
Single-channel fit & Pointlike and alternative predictions under a declared cutoff policy & Does any alternative improve held-out prediction without a nuisance explanation? & Report only a channel-conditioned preference or constraint.\\
Cross-channel synthesis & Qualified channels connected by an explicit model relation & Is a shared alternative compatible with all relevant channel likelihoods? & No-go for a common interpretation if coefficient or shape relations conflict.\\
Independent confirmation & A separate observable or facility with materially different systematics & Does the preferred pattern survive a new resolution and uncertainty environment? & Upgrade only after independent confirmation and repeated validity tests.\\
\bottomrule
\end{tabularx}
\end{table*}

\subsection{How each outcome changes the scientific answer}

The pointlike-SM branch is not an empty outcome. If it survives the qualified likelihood comparison, it supports a sharper statement: within the registered momentum-transfer interval, fiducial definition, uncertainty model, and alternative class, no additional fermion substructure is required. This statement can be compared across runs because the resolution ledger records the terms under which it is true. It remains stronger than a casual assertion that nothing was found, yet it does not outgrow its evidence.

The contact-EFT branch has a different interpretation. A stable contact coefficient can identify a departure from the pointlike baseline in an effective description, but it cannot by itself identify a constituent scale or a unique ultraviolet completion. The analysis must disclose whether the preferred coefficient is robust to changes in cutoff, dimension-eight terms, and flavor assumptions. High-energy tail studies make this issue explicit: the utility of the tail is matched by sensitivity to the domain in which the EFT description is reliable \cite{allwicher2023}.

The form-factor branch asks a still different question. A radius-like parameter can summarize a smooth momentum-transfer dependence in a valid low-$Q$ expansion. Its apparent value must be tested against response functions and PDF variations that can also tilt a spectrum. An excited-fermion branch adds resonance, acceptance, and reconstruction questions. Existing excited-lepton searches illustrate why a benchmark limit must retain its mass and coupling assumptions \cite{cms_tau2025}. RLMC keeps those hypotheses separate because their failure modes are separate.

\subsection{Strategic value of complementary facilities}

Energy, luminosity, flavor reach, and experimental systematics are not interchangeable resources. A future hadron collider can extend high-$Q$ reach; a lepton collider can emphasize cleaner initial conditions and precision; an electron-ion program can constrain partonic structure with a different kinematic organization. Their complementarity is valuable only when it is expressed through a transfer condition and status label, not through a decorative collection of sensitivity plots. Planning studies for future colliders and electron-ion measurements provide motivation for these distinct roles \cite{abada2019,arrington2021}.

The resulting scientific answer is dynamic but rigorous. The smallest currently tested fields are those whose pointlike-SM descriptions remain adequate after the best available resolution tests and stated alternatives. A future result may lower a radius bound, exclude an excited-state region, restrict an EFT coefficient, or reveal a coherent deviation. Each possibility refines the answer without requiring the proposal to claim that it has already occurred.

\appendices
\section{Provenance and Execution Contract}

The proposal inherits Survey identity \texttt{survey-elementary-matter-001}, project identity \texttt{phys\_7\_elementary\_matter\_20260901}, and the version-1 context of Standard-Model pointlike-fermion compositeness resolution. The Idea Agent selects \texttt{idea-rlmc-001}; the ExperimentDesign Agent binds it to a computational-digital design.

The execution policy is \texttt{DESIGN\_ONLY}. Observed results are an empty set, automatic execution is disabled, and human review is required. Planned synthetic closures, future-facility estimates, and public-data interfaces are marked as distinct statuses. This prevents the research plan from inventing collision observations, an exclusion scale, or a discovery claim.

\section{Claim--Evidence Matrix}

\begin{table*}[!t]
\caption{Claim--evidence matrix. All citations are frozen by the Survey Agent.}
\label{tab:evidence}
\centering
\footnotesize
\begin{tabularx}{\textwidth}{p{0.31\textwidth}p{0.27\textwidth}Y}
\toprule
\textbf{Bounded report claim} & \textbf{Evidence keys} & \textbf{Inference constraint}\\
\midrule
The SM classifies quark, lepton, gauge, and Higgs fields while leaving open scope. & \texttt{pdg2008}, \texttt{cottingham2007}, \texttt{cms\_higgs2015} & Do not transform a successful inventory into proof against all high-scale alternatives.\\
High-energy Drell--Yan tails can constrain semileptonic EFT interactions. & \texttt{allwicher2023}, \texttt{brivio2017} & Declare flavor assumptions, operator basis, EFT order, and cutoff policy.\\
Excited-fermion null searches give benchmark-dependent constraints. & \texttt{cms\_tau2025} & Do not make a global indivisibility claim from one channel or benchmark.\\
Lepton-hadron and future-collider studies are complementary. & \texttt{arrington2021}, \texttt{abada2019} & Preserve observed, synthetic, and projected status labels.\\
Low-energy and weakly coupled searches motivate a broader program. & \texttt{alekhin2016}, \texttt{jaeckel2010}, \texttt{colladay1998} & Do not present a parameterization or search motivation as detection evidence.\\
\bottomrule
\end{tabularx}
\end{table*}

For every later implementation, the record must declare the channel and data status; $Q$ definition and binning; fiducial selection; response and background model; theory calculator and PDF prescription; nuisance correlations; EFT basis, cutoff, and truncation policy; alternative-hypothesis parameters; closure-test configuration; held-out split; and decision branch from Table~\ref{tab:branches}. This is the minimum audit trail from a finite-resolution observation to a bounded conclusion.

\begin{thebibliography}{00}
\bibitem{pdg2008} C. Amsler \emph{et al.}, Review of Particle Physics, \emph{Physics Letters B}, vol. 667, no. 1--5, pp. 1--6, 2008, doi: 10.1016/j.physletb.2008.07.018.
\bibitem{cottingham2007} W. N. Cottingham and D. A. Greenwood, \emph{An Introduction to the Standard Model of Particle Physics}, 2nd ed. Cambridge, U.K.: Cambridge Univ. Press, 2007, doi: 10.1017/CBO9780511791406.
\bibitem{cms_higgs2015} V. Khachatryan \emph{et al.}, Precise determination of the mass of the Higgs boson and tests of compatibility of its couplings with the standard model predictions, \emph{Eur. Phys. J. C}, vol. 75, p. 212, 2015, doi: 10.1140/epjc/s10052-015-3351-7.
\bibitem{allwicher2023} L. Allwicher, D. A. Faroughy, F. Jaffredo, O. Sumensari, and F. Wilsch, Drell--Yan tails beyond the Standard Model, \emph{J. High Energy Phys.}, vol. 2023, no. 3, p. 064, 2023, doi: 10.1007/JHEP03(2023)064.
\bibitem{brivio2017} I. Brivio, Y. Jiang, and M. Trott, The SMEFTsim package, theory and tools, \emph{J. High Energy Phys.}, vol. 2017, no. 12, p. 070, 2017, doi: 10.1007/JHEP12(2017)070.
\bibitem{cms_tau2025} A. Hayrapetyan \emph{et al.}, Search for excited tau leptons in the tau tau gamma final state in proton-proton collisions at 13 TeV, \emph{J. High Energy Phys.}, vol. 2025, no. 6, p. 006, 2025, doi: 10.1007/JHEP06(2025)006.
\bibitem{arrington2021} J. Arrington \emph{et al.}, Revealing the structure of light pseudoscalar mesons at the electron-ion collider, \emph{J. Phys. G}, vol. 48, p. 075106, 2021, doi: 10.1088/1361-6471/abf5c3.
\bibitem{abada2019} A. Abada \emph{et al.}, FCC Physics Opportunities, \emph{Eur. Phys. J. C}, vol. 79, p. 474, 2019, doi: 10.1140/epjc/s10052-019-6904-3.
\bibitem{alekhin2016} S. Alekhin \emph{et al.}, A facility to search for hidden particles at the CERN SPS: the SHiP physics case, \emph{Rep. Prog. Phys.}, vol. 79, p. 124201, 2016, doi: 10.1088/0034-4885/79/12/124201.
\bibitem{jaeckel2010} J. Jaeckel and A. Ringwald, The Low-Energy Frontier of Particle Physics, \emph{Annu. Rev. Nucl. Part. Sci.}, vol. 60, pp. 405--437, 2010, doi: 10.1146/annurev.nucl.012809.104433.
\bibitem{colladay1998} D. Colladay and V. A. Kostelecky, Lorentz-violating extension of the standard model, \emph{Phys. Rev. D}, vol. 58, p. 116002, 1998, doi: 10.1103/PhysRevD.58.116002.
\end{thebibliography}
